% !TeX root = ../papers/nc-p3.tex
\section{三重横断線の外の三母線と四次核の基点}
\label{sec:cofactor-triple-transversal}
\begin{theorem}\label{thm:cofactor-eighteen-no-cubic}
$d=18$ で三次曲面は核像にならない。
\end{theorem}
\begin{proof}
前定理の残りは特異平面三次を底とする錐で、$S_0=\mathbb F_3$、
核束 $L=E+3F$、入力逆系 $A=E+6F$ の場合である。
自乗九なので一単純基点 $p$ を吹く。入力次数一の母線を持つため
$p$ を通る母線の固有変換を潰す初等変換で、入力射は新しい滑らかな直線束 $R$ 上の
相対超平面束 $H$ から定まり $H^2=8$ である。
入力射 $f:R\to D$ は有限である。縮約曲線があれば水平で全入力母線に共通点があり、
$D$ は錐になる。しかし整な平面八次を底とする錐の特異曲線は頂点を通る直線だけで、
十八次の $C$ がそこに入れない。

一般入力平面の核曲線 $J=G$ では頂点を核とする全ての点の入力は一つの点である。
元の負切断の固有変換は入力次数 $3$ または $2$ を持ち（後者は $p\in E$）、
$p\in E$ なら例外直線も頂点の上にあり入力次数一である。
一般平面においてこの全像が一点で切られるので像は共通の直線 $\Lambda$ である。
従って一般平面で頂点に行く正規化枝は相異なる三本である
（次数三の切断、または次数二の切断と一本の例外直線。特性零で一般値を取る）。
$J=G$ の算術種数三とその正規化の種数零より、一般入力平面上での余核から正規化への
評価の不足は全てこの三枝にある。実際その入力点は階数三で平面八次の三枝の余長は三以上である。
初等変換で潰した母線の入力点も $\Lambda$ にある。$p\notin E$ ならその母線は負切断に、
$p\in E$ なら元の例外直線に会うからである。
ゆえに $\Lambda$ 外では核束が $R$ に定まり、$M$ はその有限直像に等しい。
一般平面での等号から余次元一で等しく、$M$ の深さ二で有限集合を越えて延長するためである。
従って $\Lambda$ 外の少なくとも十七個の印は、入力母線族の孤立した長さ三繊維である。

その個数を別に計算する。入力ベクトル空間を $V$、パラメータ直線 $B=\mathbb P^1$ 上で
母線二束の商を $F'$ とする。階数二で次数八である（Plücker次数が $H^2$）。
固定直線の二空間 $W\subset V$ の像を飽和させると線束 $P\subset F'$ で、
$F'/P=Q$ も線束となる。これは全母線が $\Lambda$ に会い、一般母線はそれに一致しないからである。
$V/W$ の二切断は $Q$ を生成し、母線と $\Lambda$ の張る平面への射
$c:B\to\mathbb P^1$ を与える。$c$ は次数五である。
一般平面束の一員は $\Lambda$ の外では $c$ の一般繊維の相異なる入力母線を各一重で切る。
$\Lambda$ の一般点では重複度三である。実際上記の一般横断平面八次の三枝は合計余長三なので
三本とも滑らかで接方向は相異なる。平面束の一般平面が一般点でいずれかの接平面に一致することはない
（選んだ固定平面に一般点で一致する平面は高々三つ）。ゆえに $8=3+\deg c$。
従って $P=\mathcal O(3)$、$Q=\mathcal O(5)$ である。

$B^{(3)}=\mathbb P^3$ とし、閉曲線 $T_3\subset B^{(3)}\times\mathbb P^1$ で
長さ三の因子 $Z\le c^*(s)$ を記録する。これは二項五次式が $Z$ 上で消える三条件であり、
各 $s$ 上の部分因子の候補は有限なので局所完全交叉で底に有限平坦である。
一般の相異なる五点から三つを選ぶので底次数は十、$B^{(3)}$ の超平面次数は六である
（指定した一般一点を含む条件を課し残る四点から二点を選ぶ。どちらも一般繊維は非分岐）。
$s$ の平面の三次元ベクトル空間 $V_s$ は $W$ と $V/W$ の普遍部分直線からなり、
$\det V_s=\mathcal O(-1)$。$T_3$ 上では $V_s$ の $Q^{[3]}$ への評価は零なので
$V_s\to P^{[3]}$ の三行行列式を取れる。$P=\mathcal O(3)$ に対し
$\det P^{[3]}=\mathcal O_{B^{(3)}}(1)$ である（普遍三次因子の列で、四つの三次式を
普遍式の一本で割る）。従ってこの行列式の零因子の次数は $6+10=16$ である。

行列式は曲線成分上で恒等零にならない。$\Lambda$ 外の共通入力点では
$M$ と有限押出しの等号のため三本以上の母線の共通点は有限で、そこでは三点繊維自身が唯一の $Z$ である。
$\Lambda$ 内の一般点については、$p\notin E$ なら三母線パラメータは
次数三の写像 $B\to\Lambda$ の一繊維である。これらが一般に同じ $c$ の値を持てば
次数五の $c$ が次数三写像を経由してしまう。$p\in E$ なら一本は固定した例外母線
（その入力像は $\Lambda$ 自身）で、$s$ は固定パラメータの $c$ 値に限られる。
また有限個の特別な入力点の上でも母線パラメータ直線上の零因子の部分因子は有限である。
従って行列式は非零因子である。$\Lambda$ 外の三点繊維はこの零点を相異なるものとして与える。
一つの $Z$ に相異なる二入力点があればそれらを結ぶ直線全体が三母線条件を満たすからである。
十七個以上は次数十六に反する。
\end{proof}

\begin{proposition}\label{prop:cofactor-eighteen-quartic-no-bases}
$d=18$ の四次核曲面は錐にならない。核正規化の滑らかなモデルで固定部を除いた
逆入力系は基点なしで、自乗八、核超平面との交点六である。
\end{proposition}
\begin{proof}
まず錐では底平面四次の正規化上の母線束で $L\equiv E+4F$, $E^2=-4$。
$A\equiv aE+6F$ は移動部なので $AE\ge0$ と $A^2\ge8$ から $a=1$、自乗八で自由である。
$H\equiv E+6F$ は豊富で入力射は有限である。底種数を $g$ とすれば跡因子
$2L-K-2H\equiv2E+(2-2g)F$ は有効で、負切断の係数はちょうど二である
（$E^2=-4$ と一般母線との交点二）。
入力次数二の $E$ が円錐曲線へ同型ならその一般点は階数三、頂点の逆像は $E$ だけで、
核座標は $E$ の一般法方向に微分非零である。入力がそこで同型になり跡に反する。
従って $E$ は直線に二対一となる。定理\ref{thm:cofactor-eighteen-no-quadric}の
分裂母線の証明と同様、二滑らか枝の導手の係数二は接平面の一致を与える。
その共通平面上の核の法方向微分は両枝の核母線方向を一致させる。
しかしこれらは底の二つの一般点の平面四次への正規化像であり、双有理なので異なる。

非錐の場合も滑らかなモデルの逆移動部 $A$ について $AL=6$, $8\le A^2\le9$。
九ならHodgeの等号により $A\equiv3L/2$。ちょうど一つの単純基点 $p$ があり、
$L$ が潰す曲線上にはない（その曲線で $A$ の次数が零なら基点を持つ切断はそこ全体で消える）。
$p$ を吹いた後に $L$ 正次数の曲線が入力に潰れることもない。核超平面を核点 $y(p)$ を通して
一般に取れば、$p$ での曲線重複度はその $L$ 次数以下、よって $3L/2$ の次数を使い切れない。
例外 $e$ の入力次数は一、その入力直線の一般点は階数三である。
$y(p)$ の逆像に他の点や例外曲線があっても、それらが動く入力を与えることはない
（他の点では基点なしで、$L$ 次数零の全曲線の入力も一点）。
核射の $p$ での微分は非零である。零なら $y(p)$ を通る二つの一般超平面の引戻しが
$p$ に少なくとも四の交点を持つ。一方核が非錐なら $y(p)$ の重複度は四未満で、
二平面の共通直線はその外でも曲面を切る。共通固定例外曲線があれば先に除くと、
移動部の自乗は $L^2$ 以下である（$L$ と直交する固定例外類の自乗が非正）のでやはり矛盾する。

一般の $e$ の点では入力は $e$ に沿って非分岐、核射は法方向に非零である。
従って入力有限な開集合で、階数三の入力射はその一般点で同型である。
しかし跡因子の類 $2L-K-2H$ は $e$ との交点が $-1$ である。
極を許す入力縮約曲線は元の $L$ 次数零の曲線だけで $p$ と交わらず、$e$ とも交わらない。
従って $e$ は零因子でなければならず、同型点の跡に反する。
\end{proof}
\path{computations/verify_cofactor_triple_transversal.py} は三母線行列式の二つの次数を検算する。
